\documentclass[a4paper,11pt]{article}
\usepackage[utf8]{inputenc}
\usepackage[T1]{fontenc}
\usepackage[french]{babel}
\usepackage[left=2cm,right=2cm,top=2cm,bottom=2cm]{geometry}
\usepackage{xcolor}
\usepackage{hyperref}
\usepackage{pifont}
\usepackage{parskip}

% --- Couleurs Crédit Agricole (vert) ---
\definecolor{primary}{RGB}{0, 100, 62}
\hypersetup{colorlinks=true, urlcolor=primary}

\begin{document}

% --- EXPÉDITEUR ---
\noindent
\begin{minipage}[t]{0.5\textwidth}
    \textbf{Loïc Aron MBASSI EWOLO}\\
    Élève-Ingénieur Bac+5 -- Double diplôme ENSEA\\[3pt]
    Cergy, France\\
    (+33) 7 59 52 33 98\\
    \href{mailto:loicaron.mbassiewolo@ensea.fr}{loicaron.mbassiewolo@ensea.fr}
\end{minipage}
\hfill
% --- DATE ---
\begin{minipage}[t]{0.4\textwidth}
    \raggedleft
    Cergy, le 20 mars 2026
\end{minipage}

\vspace{1.2cm}

% --- DESTINATAIRE ---
\hfill
\begin{minipage}[t]{0.5\textwidth}
    \raggedleft
    \textbf{Crédit Agricole Technologies et Services}\\
    Service Recrutement\\
    Paris, France
\end{minipage}

\vspace{1cm}

\noindent\textbf{Objet :} Candidature au poste de Développeur Back-End en alternance (Tribu Marketing Digital)

\vspace{0.8cm}

Madame, Monsieur,

Concevoir une API REST capable de gérer des transactions bancaires concurrentes en toute sécurité, puis architecturer une plateforme de messagerie réactive traitant des milliers de connexions simultanées via Spring Boot WebFlux : ces défis techniques ont structuré mes trois dernières années de développement backend. C'est avec un intérêt particulier que je postule à l'alternance Développeur Back-End au sein de la squad CRM Opérationnel de CA-TS, où l'alimentation des leads via Kafka et le développement d'API REST Java rejoignent directement mon parcours.

Lors de ma mission chez CSC, j'ai conçu le backend d'une plateforme de transactions financières sécurisée dans le secteur bancaire. Ce projet m'a confronté aux exigences de conformité propres aux systèmes financiers, à la gestion de la concurrence des accès aux données et à la mise en place d'une architecture scalable (Docker, Redis, tests unitaires). En parallèle, j'ai maintenu pendant deux ans la plateforme de télémédecine SOSDOCTA, où la gestion de données sensibles (RGPD) et l'authentification sécurisée (OAuth2/JWT) étaient au c\oe ur du développement.

Sur le plan de l'architecture distribuée, j'ai participé au projet Snappy, une plateforme de messagerie d'entreprise construite sur Spring Boot WebFlux avec programmation réactive (Mono/Flux), Redis Pub-Sub pour la distribution des messages, et PostgreSQL avec indexation composite optimisée. J'ai également conçu une architecture microservices complète avec GoFind, intégrant RabbitMQ pour la communication asynchrone entre services, une expérience directement transposable à l'écosystème Kafka de CA-TS. Mon stage en cours à l'INRIA Saclay (équipe TRIBE, plateau de Saclay) renforce ma rigueur méthodologique à travers le développement de pipelines Python d'analyse de données à grande échelle.

Je suis disponible dès septembre 2026 pour une alternance de 12 mois, à l'issue de mon stage INRIA. Travailler au sein d'une équipe agile pluridisciplinaire de 12 personnes, accompagné par un tuteur dédié, correspond à l'environnement dans lequel je donne le meilleur de moi-même. La perspective de contribuer aux outils qui impactent le quotidien de millions de clients bancaires du réseau Crédit Agricole est une motivation forte.

Je reste à votre disposition pour un entretien, notamment pour l'exercice de code mentionné dans le processus de recrutement.

Je vous prie d'agréer, Madame, Monsieur, l'expression de mes salutations distinguées.

\vspace{0.8cm}

\textbf{Loïc Aron MBASSI EWOLO}\\
\textit{Élève-Ingénieur ENSEA / Polytechnique Yaoundé}

\end{document}
